\scanpage{16}\mspage{15}
% ink top, heading of the three equalities
trois composantes (relatives à la \unsure{donné|donnée} de $\{a,b,c\}$)
\[
\begin{aligned}
\deleted{\unsure{C}}X_{\{a,b\}}&=\deleted{\unsure{C(a,b,c)}}=(\text{comp. de }C(\{a,b\})\text{ contenant }c)=\mathcal L_{>b}\\
X_{\{b,c\}}\deleted{\unsure{C(b,c,a)}}&=(\rule{1.45cm}{.4pt}\ C(\{b,c\})\ \rule{1.1cm}{.4pt}\ a)=\mathcal L_{<b}\\
X_{\{c,a\}}\deleted{C(c,a,b)}&=\rule{1.2cm}{.4pt}\ C(\{c,a\})\ \rule{.85cm}{.4pt}\ b)=\deleted{\unsure{\mathcal L}}\mathopen{]}a,c\mathclose{[}
\end{aligned}
\]

% ink middle, On trouve and three intersections
On trouve
\[
\begin{aligned}
&\deleted{C}X_{\{a,b\}\deleted{\unsure{c}}}\cap\deleted{C}X_{\{b,c\}\deleted{\unsure{a}}}=\varnothing\\
&\left.\begin{aligned}
&\deleted{C}X_{\{a,b\}\deleted{\unsure{c}}}\cap\deleted{C}X_{\{c,\deleted{\unsure{b}}a\}\deleted{\unsure{b}}}=\mathopen{]}b,c\mathclose{[}\ne\varnothing\\
&X_{(b,\unsure{c|b}\}}\cap X_{\{c,a\}}=\mathopen{]}a,b\mathclose{[}\ne\varnothing
\end{aligned}\right\}
\end{aligned}
\]
% ink right of the two nonempty intersections
\margin{divisibilité !}
% ink lower middle, donc b
donc $b$ est distingué, dans $J=\{a,b,c\}$, par la propriété que pour les deux \deleted{parties} ${\unsure{\varepsilon|\mathcal E|E}}',\unsure{\varepsilon|\mathcal E|E}'$ \deleted{de} $\in\unsure{P|\mathscr P|\mathfrak P}_2(J)$ qui le \unsure{contient|contiennent}, on a
\[
 X_{\unsure{\varepsilon_{0}|\varepsilon_{,}|\varepsilon}}\wedge X_{\unsure{\varepsilon|\mathcal E|E}'}=\varnothing\,.\qquad\text{On gagne !}
\]

% ink bottom, NB
\uline{NB} On a à nouveau utilisé la divisibilité.

% pale vertical watermark
\rotatebox{90}{\textsc{\unsure{ng Extra Strong}}}
